\documentclass{article}
\usepackage{PRIMEarxiv} % Core package for the PrimeAI Template
\usepackage{amsmath, amssymb, amsthm, bm, dcolumn} % Add amsthm here for the proof environment
\usepackage[numbers,sort&compress]{natbib} % Natbib for citations
\usepackage{graphicx} % For high-quality images
\usepackage{multicol} % Multi-column figures/tables
\usepackage{hyperref} % Hyperlinks
\usepackage{listings} % Code listings
\usepackage{authblk} % For structured affiliations
% Define theorem style (optional)
\theoremstyle{plain}
\newtheorem{theorem}{Theorem}
\renewcommand{\qedsymbol}{}

% Custom commands for primes and qubit encoding
\newcommand{\primeQubit}[1]{|p_{#1}\rangle}
\newcommand{\primeGate}[1]{U_{p_{#1}}}

\usepackage{wrapfig}
\usepackage[pscoord]{eso-pic}
\usepackage[fulladjust]{marginnote}
\reversemarginpar

% Typesetting improvements without footnote patching
\usepackage[protrusion=true, expansion=true, tracking=false]{microtype}
\microtypecontext{spacing=nonfrench}

% Line numbers
\usepackage[right]{lineno}

% Text layout - adjust as needed
\raggedright
\setlength{\parindent}{0.5cm}
\textwidth 5.25in 
\textheight 8.75in

% Set double spacing
\usepackage{setspace} 
\doublespacing

% Adjust width for specific content
\usepackage{changepage}

% Adjust caption style
\usepackage[aboveskip=1pt,labelfont=bf,labelsep=period,singlelinecheck=off]{caption}

% Remove brackets from references
\makeatletter
\renewcommand{\@biblabel}[1]{\quad#1.}
\makeatother

% Header, footer, and page numbers
\usepackage{lastpage,fancyhdr}
\pagestyle{fancy}
\fancyhf{}
\fancyhead[L]{Preprint - PrimeAI Enhanced Template}
\fancyfoot[C]{\scriptsize Multiplicity Theory © 2025 Citizen Gardens \\ Licensed Under MIT and CC BY-NC-SA 4.0.}
\fancyfoot[R]{Page \thepage\ of \pageref{LastPage}}
\renewcommand{\footrule}{\hrule height 2pt \vspace{2mm}}

\begin{document}

\title{Universal Multiplicity Equation: \\ A Next-Generation Framework}

\author[1]{Ryan O. Van Gelder}
\author[1]{Nicholas Galioto}
\author[2]{Chris McGinty}
\affil[1]{Citizen Gardens}
\affil[2]{Skywise AI}
\date{\today}

\maketitle
\begin{abstract}
The Enhanced Universal Multiplicity Equation (EUME) is presented as an advanced mathematical framework that integrates multi-scale dynamics, holographic principles, and quantum-inspired geometries. This work builds on the foundational Universal Multiplicity Equation (UME) by incorporating elements from the McGinty Equation (MEQ), which encapsulates golden ratios, sacred geometry, quantum fluctuations, and cosmological horizons. The EUME is designed to model complex, interconnected systems spanning quantum, cosmological, and classical domains.
\end{abstract}
\section{The Universal Multiplicity Equation}
The general form of the Universal Multiplicity Equation is expressed as:
\begin{equation}
\frac{\partial \psi_k(t)}{\partial t} = \alpha_k(t)\psi_k + \beta_k(t) \int_0^t I_k(\tau) d\tau + \gamma_k(t) \sum_{j,l} T_{kjl} \psi_j \psi_l + \lambda_k(t) \nabla^2 \psi_k + \eta_k(t) \psi_k^n + \xi_k(t),
\end{equation}
where:
\begin{itemize}
    \item $\psi_k(t)$: The state variable evolving over time.
    \item $\alpha_k(t)$: Growth or decay coefficient, dynamically adjustable.
    \item $\beta_k(t)$: Memory coefficient for historical integration.
    \item $I_k(\tau)$: Input function capturing external influences.
    \item $\gamma_k(t)$: Coupling coefficient for tensor-driven interactions.
    \item $T_{kjl}$: Higher-order tensor encoding multi-scale dependencies.
    \item $\lambda_k(t)$: Quantum potential term for wave-like behavior.
    \item $\eta_k(t)$: Nonlinear self-interaction coefficient.
    \item $n$: Degree of nonlinearity.
    \item $\xi_k(t)$: Stochastic noise term modeling randomness.
\end{itemize}

\subsection{Tensor-Driven Multi-Scale Interactions}
Higher-order tensor interactions are modeled as:
\begin{equation}
\sum_{j,l} T_{kjl} \psi_j \psi_l,
\end{equation}
which captures dependencies across multiple dimensions, enabling scalable and interconnected computations.

\subsection{Quantum Potential and Wave Propagation}
The quantum potential term introduces spatial coherence and wave dynamics:
\begin{equation}
\lambda_k(t) \nabla^2 \psi_k,
\end{equation}
reflecting quantum-inspired properties in evolving systems.

\subsection{Nonlinearity and Emergence}
The nonlinear self-interaction term:
\begin{equation}
\eta_k(t) \psi_k^n,
\end{equation}
accounts for emergent behaviors such as saturation, bifurcation, and chaotic dynamics.

\subsection{Dynamic Memory and Feedback}
Historical effects are incorporated via:
\begin{equation}
\beta_k(t) \int_0^t I_k(\tau) d\tau,
\end{equation}
allowing the system to learn and adapt based on past states.

\subsection{Stochasticity}
Environmental randomness is modeled through a stochastic noise term:
\begin{equation}
\xi_k(t) \sim \mathcal{N}(0, \sigma_k^2),
\end{equation}
representing uncertainties and perturbations.

\subsection{Expanded Representation}
To unify various scales and domains, the UME can be represented in tensor form:
\begin{equation}
\frac{\partial \bm{\Psi}(t)}{\partial t} = \mathcal{A}(t) \bm{\Psi}(t) + \mathcal{B}(t) \int_0^t \bm{I}(\tau) d\tau + \mathcal{T}(t) \otimes \bm{\Psi}(t) \otimes \bm{\Psi}(t) + \mathcal{Q}(t) \nabla^2 \bm{\Psi}(t) + \mathcal{E}(t),
\end{equation}
where $\mathcal{A}(t)$, $\mathcal{B}(t)$, and $\mathcal{T}(t)$ are dynamic tensors, and $\mathcal{E}(t)$ encodes stochastic noise.

\section{Enhanced Universal Multiplicity Equation}
The Enhanced Universal Multiplicity Equation (EUME) provides a unified representation for modeling interactions across scales, integrating:
\begin{itemize}
    \item Sacred geometry (\(\mathcal{G}\)) and the golden ratio (\(\Phi\)).
    \item Energetic variations (\(\Gamma\)) shaping geometric configurations.
    \item Cosmological horizons (\(\Lambda\)) affecting quantum and classical dynamics.
    \item Quantum fluctuations (\(\mathcal{Q}\)) as drivers of stochastic processes.
    \item Fundamental constants (\(\epsilon_0, \mu_0, \pi\)) linking physical and mathematical domains.
\end{itemize}


The fully enhanced EUME is expressed as:
\begin{equation}
\frac{\partial \psi_k(t)}{\partial t} = \alpha_k(t)\psi_k + \beta_k(t) \int_0^t I_k(\tau, \Phi, \Gamma) \, d\tau + \gamma_k(t) \sum_{j,l} T_{kjl}(\Phi, \mathcal{G}, \epsilon_0, \mu_0) \psi_j \psi_l + \Lambda_k(t, \Phi, \Gamma)\nabla^2 \psi_k + \xi_k(t, \Phi) + \int_{\Omega} H_{McGinty} \, d\Omega.
\label{eq:EUME}
\end{equation}

\section{Breakdown of Components}
\subsection{Core Terms of EUME}
\begin{enumerate}
    \item \textbf{State Evolution:}
    \[
    \frac{\partial \psi_k(t)}{\partial t}
    \]
    Represents the temporal evolution of the state variable \(\psi_k(t)\).

    \item \textbf{Growth/Decay Term:}
    \[
    \alpha_k(t)\psi_k
    \]
    Encodes dynamic growth or decay behavior modulated by \(\alpha_k(t)\).

    \item \textbf{Memory and Feedback:}
    \[
    \beta_k(t) \int_0^t I_k(\tau, \Phi, \Gamma) \, d\tau
    \]
    Captures historical effects scaled by energetic variations \(\Gamma\) and influenced by sacred geometry \(\Phi\).

    \item \textbf{Tensorial Interactions:}
    \[
    \gamma_k(t) \sum_{j,l} T_{kjl}(\Phi, \mathcal{G}, \epsilon_0, \mu_0) \psi_j \psi_l
    \]
    Describes multi-dimensional interactions, where \(T_{kjl}\) incorporates geometric and physical constants.

    \item \textbf{Quantum Potential Term:}
    \[
    \Lambda_k(t, \Phi, \Gamma)\nabla^2 \psi_k
    \]
    Models spatial coherence and wave propagation, linked to cosmological expansion.

    \item \textbf{Stochastic Noise:}
    \[
    \xi_k(t, \Phi) = \mathcal{Q}(\Phi \cdot \Gamma) + \epsilon \cdot \sin(\omega t)
    \]
    Introduces randomness due to quantum fluctuations and harmonic oscillations.

    \item \textbf{Information Integral:}
    \[
    \int_{\Omega} H_{McGinty} \, d\Omega
    \]
    Represents the integration of encoded information across the domain \(\Omega\), adhering to the holographic principle.
\end{enumerate}

\subsection{McGinty Equation Integration}
The McGinty term is defined as:
\[
H_{McGinty} = \int \left( \Sigma(\Phi \otimes \mathcal{G}) \cdot \Lambda \cdot \frac{\partial}{\partial \mathcal{Q}} \left(\Phi \cdot \sqrt{\Gamma}\right) \right) d\Omega,
\]
where:
\begin{itemize}
    \item \(\Phi\): Golden ratio, representing harmonic proportionality.
    \item \(\mathcal{G}\): Sacred geometry, encoding fractal-like symmetries.
    \item \(\Lambda\): Cosmological horizon, linked to universe expansion.
    \item \(\mathcal{Q}\): Quantum vacuum, describing fluctuations and zero-point energy.
    \item \(\Gamma\): Energetic variation, coupling energy to geometry.
\end{itemize}

\subsection{Conclusion}
The EUME represents a comprehensive framework that unifies geometry, energy, and quantum dynamics. Future work will focus on numerical simulations and domain-specific applications, including quantum computing, cosmology, and artificial intelligence.

\section{Prime-Encoded EUME}
The prime-encoded Enhanced Universal Multiplicity Equation (EUME) is expressed as:
\begin{equation}
p_\psi^k \cdot p_t = p_\alpha^{k+t} \cdot p_\psi^k 
+ p_\beta^{k+t} \cdot p_\int \prod_{p_\Phi, p_\Gamma} p_t 
+ p_\gamma^{k+t} \sum_{j,l} \left( \prod_{p_j, p_l} p_T^{j \cdot l} \cdot p_\psi^j \cdot p_\psi^l \right)
+ p_\Lambda^{k+t} \cdot p_\nabla^2 \cdot p_\psi^k 
+ p_\xi^{k+t} 
+ p_\int \prod_{p_\Phi, p_\mathcal{G}, p_\Gamma, p_\mathcal{Q}, p_\Lambda}.
\label{eq:prime-eume}
\end{equation}

\subsection{Component Breakdown}
\begin{enumerate}
    \item \textbf{State Variable:}
    \[
    \psi_k(t) \rightarrow p_\psi^k = p_k^t
    \]

    \item \textbf{Coefficients:}
    \[
    \alpha_k(t) \rightarrow p_\alpha^{k+t}, \quad 
    \beta_k(t) \rightarrow p_\beta^{k+t}, \quad
    \gamma_k(t) \rightarrow p_\gamma^{k+t}, \quad
    \Lambda_k(t) \rightarrow p_\Lambda^{k+t}.
    \]

    \item \textbf{Tensor Interactions:}
    \[
    T_{kjl} \rightarrow \prod_{p_j, p_l} p_T^{j \cdot l}.
    \]

    \item \textbf{Stochastic Noise:}
    \[
    \xi_k(t) \rightarrow p_\xi^{k+t}.
    \]

    \item \textbf{McGinty Integral:}
    \[
    H_{McGinty} \rightarrow \prod_{p_\Phi, p_\mathcal{G}, p_\Gamma, p_\mathcal{Q}, p_\Lambda}.
    \]
\end{enumerate}

\sectionsub{Implications of Prime Encoding}
\begin{itemize}
    \item \textbf{Efficiency:} Prime encoding ensures modular representations that can be efficiently manipulated in symbolic and computational frameworks.
    \item \textbf{Quantum-Inspired Dynamics:} By leveraging prime-based modular arithmetic, the equation supports quantum-inspired optimization and error correction.
    \item \textbf{Scalability:} The encoding is scalable across high-dimensional systems, bridging micro and macro scales.
\end{itemize}

\subsection*{Conclusion}
Prime encoding extends the Enhanced Universal Multiplicity Equation (EUME) by embedding modular and symbolic computational capabilities, enhancing its utility for quantum systems, cryptography, and AI.

\nocite{*}
\bibliographystyle{unsrt}
\bibliography{references}

\end{document}